\chapter{Introducción}
\label{chap:Introduction}

\section{Contexto}

La adopción de redes \textit{Internet of Things} (IoT) e \textit{Industrial IoT} (IIoT) ha aumentado la superficie de ataque de hogares, laboratorios, plantas industriales y sistemas de monitoreo conectados. Estos entornos combinan dispositivos heterogéneos, protocolos livianos como MQTT, conectividad intermitente y capacidades limitadas de cómputo, lo que dificulta desplegar mecanismos de detección convencionales basados en procesamiento centralizado. Bajo marcos de gestión de riesgo como el \textit{NIST Cybersecurity Framework}, la detección temprana y el monitoreo continuo son capacidades fundamentales para reducir el tiempo de respuesta ante incidentes \cite{NISTCSF2024}.

Los sistemas de detección de intrusiones de red (NIDS) basados en aprendizaje automático han mostrado altos niveles de precisión en condiciones de laboratorio. Sin embargo, muchos de estos resultados se obtienen bajo supuestos que no siempre se cumplen en IoT: centralización de datos, disponibilidad de servidores con alta memoria, ausencia de cifrado en las comunicaciones internas del sistema y modelos demasiado grandes para ejecutarse en microcontroladores. En consecuencia, el reto no consiste únicamente en maximizar \textit{accuracy}, sino en validar una arquitectura completa que pueda operar sobre hardware restringido, actualizarse de manera distribuida y proteger los artefactos del aprendizaje.

Esta tesis aborda ese reto mediante una arquitectura Edge--Fog--Cloud para detección de intrusiones en tráfico MQTT. La capa Edge se implementa con nodos ESP32-S3 que ejecutan inferencia TinyML y transmiten características de flujo; la capa Fog se implementa sobre Raspberry Pi 4 como gateway MQTT, entrenador local y agregador intermedio; y la capa Cloud se implementa en un PC con FastAPI para coordinar la agregación global. El sistema se evaluó en varias configuraciones de la versión \texttt{hfl\_v7}: red neuronal MLP con ASCON-128, red neuronal MLP sin ASCON como línea base, variante CNN-1D y estrategia FOG con agregación entre gateways.

\section{Problemática}

La primera brecha identificada es la \textbf{viabilidad de despliegue}. Modelos centralizados como Random Forest o SVM-RBF pueden alcanzar precisiones cercanas al 100\%, pero sus requerimientos de memoria son incompatibles con microcontroladores. En contraste, un IDS IoT real necesita inferencia frecuente, bajo consumo de memoria y comunicación limitada.

La segunda brecha corresponde a la \textbf{privacidad y distribución de datos}. Centralizar capturas PCAP o flujos completos desde múltiples redes introduce costos de transferencia y riesgos de exposición. El aprendizaje federado reduce este problema al compartir pesos o gradientes en lugar de datos crudos, pero en IoT introduce retos adicionales: datos \textit{non-IID}, rondas incompletas, heterogeneidad de hardware y sobrecarga de comunicación \cite{Imteaj2022_FL_ResourceConstrainedSurvey, Albanbay2025}.

La tercera brecha es la \textbf{seguridad del propio ciclo federado}. Aunque FL evita mover datos crudos, los pesos del modelo, métricas locales y mensajes de control siguen siendo activos sensibles. Por tanto, las comunicaciones entre ESP32, Raspberry Pi y PC deben proteger confidencialidad e integridad con mecanismos compatibles con dispositivos restringidos. ASCON-128, estandarizado por NIST como criptografía ligera, es una opción pertinente para este escenario \cite{NISTSP8002322025}.

Finalmente, el sistema debe ser evaluado como arquitectura, no solo como modelo. Por eso se comparan cuatro dimensiones: desempeño del modelo TinyML, costo del aprendizaje federado, overhead de ASCON-128 y efecto de variar el modelo o la topología mediante las ramas \texttt{hfl\_v7-RN}, \texttt{hfl\_v7-no-ascon} y \texttt{hfl\_v7-CNN}.

\section{Justificación}

La solución propuesta aporta una validación práctica de un IDS federado para IoT sobre hardware real. A diferencia de un entrenamiento centralizado, el sistema opera con un modelo compacto de 13 características de flujo y tres clases de tráfico: \texttt{normal}, \texttt{mqtt\_bruteforce} y \texttt{scan\_A}. Este alcance permite instrumentar el ciclo completo de detección, entrenamiento local, agregación global y redistribución del modelo sin depender de infraestructura de alto costo en el borde.

El proyecto también justifica la comparación entre variantes. La rama con ASCON-128 permite medir el costo de proteger todos los canales; la rama sin ASCON sirve como línea base para aislar el overhead criptográfico; la variante CNN-1D evalúa si una arquitectura más expresiva mejora el desempeño sin romper la viabilidad edge; y el modo FOG introduce una agregación intermedia entre Raspberry Pi antes del servidor central.

El componente de NLP sobre logs, considerado en objetivos tempranos del proyecto, no forma parte del despliegue final de \texttt{hfl\_v7}. En esta versión se deja como línea futura, mientras que el aporte implementado se concentra en características numéricas de flujo, TinyML, HFL y cifrado ligero. Esta delimitación es importante para mantener consistencia entre código, resultados y conclusiones.

El documento está estructurado de la siguiente manera: el Capítulo~\ref{chap:Introduction} presenta el problema y los objetivos; el Capítulo~\ref{chap:theo} resume el marco conceptual y el estado del arte; el Capítulo~\ref{chap:meth} describe la arquitectura y el protocolo experimental; el Capítulo~\ref{chap:resu} presenta resultados; el Capítulo~\ref{chap:conc} concluye y plantea trabajos futuros; y los anexos resumen archivos y artefactos del sistema.

\section{Objetivos}

\subsection{Objetivo General}
Diseñar, implementar y evaluar un sistema de detección de intrusiones para tráfico IoT/MQTT basado en TinyML, aprendizaje federado jerárquico y cifrado ligero ASCON-128, desplegado sobre una arquitectura Edge--Fog--Cloud compuesta por ESP32-S3, Raspberry Pi 4 y un servidor PC.

\subsection{Objetivos Específicos}
\begin{enumerate}
    \item Construir un \textit{pipeline} de datos basado en flujos de red unidireccionales, seleccionando 13 características computables en dispositivos de borde y definiendo un problema triclase para tráfico \texttt{normal}, \texttt{mqtt\_bruteforce} y \texttt{scan\_A}.

    \item Entrenar y comparar modelos compactos para IDS IoT, incluyendo MLP, CNN-1D, Residual-MLP y Tiny Transformer, exportando pesos y parámetros de normalización para ejecución en ESP32-S3.

    \item Implementar una arquitectura HFL bidireccional donde los gateways Raspberry Pi entrenan localmente, envían pesos al servidor central, reciben el modelo global y lo redistribuyen a los nodos ESP32.

    \item Evaluar el impacto del cifrado ASCON-128 frente a una línea base sin cifrado, midiendo tiempo de cifrado/descifrado, tamaño de mensajes y efecto sobre la convergencia del modelo.

    \item Analizar el efecto de variar la arquitectura del modelo y la topología de agregación mediante las ramas \texttt{hfl\_v7-RN}, \texttt{hfl\_v7-no-ascon}, \texttt{hfl\_v7-CNN} y el modo FOG.
\end{enumerate}

\subsection{Preguntas de investigación}
\begin{enumerate}
    \item ¿Puede un modelo TinyML ejecutado en ESP32-S3 detectar tráfico MQTT malicioso con precisión operativamente viable?
    \item ¿Qué reducción de transferencia de datos se obtiene al usar HFL en lugar de centralizar los flujos de red?
    \item ¿Cuál es el costo temporal y de comunicación de proteger el ciclo federado con ASCON-128?
    \item ¿Cómo cambian los resultados al comparar MLP, CNN-1D, baseline sin ASCON y agregación FOG?
\end{enumerate}
